% ==============================================================================
% T0 Theory / FFGFT: Document 272 (Chapter Content, EN) -- revised
% Title: FFGFT (T0) and HLV, classified through Dot Theory (governance version)
% Author: Johann Pascher
% Date: June 2026
% References: Doc. 190, 206, 270, 271; Dot Theory (Vossen) + IPI governance; HLV (Krueger)
% ==============================================================================

\section*{Abstract}
This work classifies FFGFT (T0, time--mass duality) and the Helix--Light--Vortex programme
(HLV, M.~Krüger) through the \emph{current governance scaffolding} of Stefaan Vossen's \emph{Dot Theory},
as developed in the direct IPI exchange (Vossen with J.~Pascher and others). Dot Theory is not a competing
physical model but a representational meta-programme with a governance machinery of four record objects:
the \textbf{Lexicon} (structure/terminology), the \textbf{OAP} (admissibility record), the \textbf{Matrix}
(relationships between frameworks), and the \textbf{FAH} (\emph{Framework Admissibility History}: how
admissibility changes; provenance-preserving renderings $O0 \to O1/O2$). Guiding principles:
\enquote{declared provenance, not universal agreement} and \enquote{preservation of meaning-making, not
allocation of meaning} -- an FAH record is a \emph{historical object}, not a truth claim. The present
classification is produced as a \textbf{Matrix entry} (relationship FFGFT$\leftrightarrow$HLV) and offered
as a \textbf{secondary rendering ($O2$)} under provenance -- a \emph{declared} reading (J.~Pascher), while
Krüger (HLV) and Vossen (Dot) hold the $O0$ originator declaration. Result: FFGFT
(observer-independent/periodic/forward-predictive) and HLV (observer-embedded/aperiodic/substrate-distinctive)
are commensurable only on the spectral/recovery axis; the other axes lie in different \enquote{question
geometries}. Only Dot's governance/routing layer is used, not its own physical ToE/$\psi$ layer.

% ============================================================================
\section{Dot Theory's governance machinery (state IPI, June 2026)}
% ============================================================================

Vossen explicitly does not present Dot Theory as a completed physical theory, but as a representational
hypothesis about modelling, with a strict programme requirement: any extension must be consistent, preserve
required symmetries, \emph{recover} the standard formalism as a limiting case, and yield at least one clear
empirical discriminator. In the direct IPI exchange this has grown into a governance machinery with four
record objects:

\begin{itemize}
\item \textbf{Lexicon} -- holds the \emph{structure}: terms and their definitions.
\item \textbf{OAP} -- holds the \emph{admissibility} of a framework (does it meet the programme requirement?).
\item \textbf{Matrix} -- holds the \emph{relationships} between frameworks.
\item \textbf{FAH} (Framework Admissibility History) -- holds \emph{how} admissibility changes when an
implicit commitment is exposed and relocated.
\end{itemize}

The FAH works with staggered renderings: $O0$ (originator declaration) $\to$ $O1/O2$ (secondary renderings
by other observers) $\to$ later evaluation. Once admitted, records remain \emph{permanent} and contestable;
later entries may refine, contest, supplement, or supersede, but the original rendering remains part of the
developmental history. The load-bearing distinction is \enquote{preservation of meaning-making} (admissible:
recording \emph{how} a change was understood) versus \enquote{allocation of meaning} (inadmissible: deciding
\emph{who} is right). In Pascher's words: \emph{the history of how a result was understood is itself part of
the result.}

This classification uses only the governance/routing layer. Dot Theory's own physical ToE layer
(meta-Lagrangian $\mathcal{L}_{\text{Dot}}$, observer operator $\odot$, observer-purpose field $\psi$) is
\emph{not} invoked. The state is June~2026 and -- consistent with Dot's non-absolutist stance and FAH
supersedability -- explicitly provisional.

% ============================================================================
\section{Lexicon: the classification axes as defined terms}
% ============================================================================

\begin{itemize}
\item \emph{Informational accessibility / question geometry}: under what question and with what access does
the framework operate?
\item \emph{Operational invariants}: what is preserved under the characteristic operations?
\item \emph{Recoverability / projection-loss}: what information is lost in projection, what can be recovered?
\item \emph{Observer embedding}: is an observer constitutively part of the formalism?
\item \emph{Empirical discriminator (admissibility)}: which testable observable separates the framework from
null models?
\end{itemize}

% ============================================================================
\section{OAP: admissibility of FFGFT and HLV}
% ============================================================================

Both frameworks are \emph{admissible} in the OAP sense, but on different observables. \textbf{FFGFT} meets the
requirement via dimensionless ratios ($\alpha$, CMB peak ratios from $D=4$), an $O(\xi)$ recovery of continuum
dispersion, and a \emph{worked-out} recoverability typology (Doc.~270). \textbf{HLV} meets it via the
graph-Laplacian spectrum against null models and vortex benchmarks; recovery of QFT/GR in the limit is required
but yet to be shown -- a \emph{maturity difference}, not a contradiction.

% ============================================================================
\section{Matrix: the relationship FFGFT $\leftrightarrow$ HLV}
% ============================================================================

\begin{center}
\small
\begin{tabular}{@{}p{0.17\linewidth}p{0.39\linewidth}p{0.39\linewidth}@{}}
\hline
\textbf{Lexicon axis} & \textbf{FFGFT / T0} & \textbf{HLV}\\
Structural nesting (June 2026)
& crystallographic $3$-fold core $C_3$; $T^4$
& icosahedral extension $A_5\supset C_3$; $T^7=T^4{\times}T^3$; $\varphi$ = eigenvalue of the perp sector\\
\hline
Accessibility / question geometry
& observer-\emph{independent}; \enquote{which dimensionless ratios follow from $\xi$?}
& substrate-centred; \enquote{is the substrate distinguishable from null models?}\\
Operational invariants
& $\tilde T\cdot m=1$; $\xi$; winding numbers; $Z_3$ ($\det=1-\xi$)
& cut-and-project window; vortex topology; triadic time $\psi=t+i\varphi+j\chi$\\
Recoverability / projection-loss
& \emph{worked out} (Doc.~270): lossy reduction vs.\ lossless representation
& recovery of QFT/GR in the limit -- required, yet to be shown\\
Observer embedding
& none (purely geometric)
& yes: $\psi$/consciousness layer (like Dot Theory itself)\\
Empirical discriminator
& dimensionless ratios: $\alpha$, peak ratios from $D=4$
& graph-Laplacian spectrum vs.\ null model; vortex benchmarks\\
\hline
\end{tabular}
\end{center}

Axis by axis: FFGFT and HLV pose \emph{different} questions (what follows from $\xi$ vs.\ is the substrate
distinguishable), carry the triadic structure at different places (time algebra $\psi$ vs.\ reduction symmetry
$Z_3$ -- an isomorphism candidate, not a proof; cf.\ Doc.~271), and connect only through the
recoverability/discriminator axis, \emph{spectrally}. There -- and only there -- is the pairwise comparison
(Doc.~271) licensed and yields the distinguishing prediction: periodic $T^4$ harmonic spectrum (peak ratios
from $D=4$) versus aperiodic quasicrystal gap spectrum. On the axes of question geometry and observer embedding
the two live in different regimes; a direct \enquote{X is like Y} comparison there would be a category error.

\medskip\noindent\textbf{Refinement (later $O2$ rendering, 19 June 2026).} On the structural / group-theoretic axis a closer relation holds than recorded above: $C_3<A_5$ and $T^4\subset T^7$ ($7=4+3$) give a \emph{containment} HLV~$\supset$~FFGFT (Doc.~283/285) -- FFGFT as the crystallographic $3$-fold core, HLV as the icosahedral extension. The golden ratio $\varphi$ is the HLV-side \emph{discriminator}: the eigenvalue of the perp directions (golden inflation $\varphi^2=\varphi+1$) that FFGFT does not contain; the same crystallographic restriction ($2\cos(2\pi/5)=1/\varphi\notin\mathbb{Z}$) that forbids $\varphi$ in FFGFT assigns it to HLV. This reading is \emph{more committed} than the original entry and is entered explicitly as a contestable rendering: it does not change OAP and allocates no meaning; the group-theoretic fact ($C_3<A_5$, $7=4+3$) is solid, the physical identification of the perp sector remains a declared reading.

% ============================================================================
\section{FAH: rendering status of this classification}
% ============================================================================

This Matrix entry is entered into provenance as a \emph{secondary rendering} ($O2$). The $O0$ originator
declarations rest with Krüger (HLV) and Vossen (Dot Theory); the reading presented here is J.~Pascher's and does
not claim to determine objectively which framework is \enquote{correct}. It \emph{preserves meaning-making} (it
records how the relationship is understood) and \emph{allocates no meaning} (it adjudicates no precedence). The
classification thus also explains \emph{retroactively} why Doc.~271 needed exactly its caveats: these are not
politeness but the consequence of different accessibility conditions -- Dot's \enquote{epistemic routing prior to
mechanism deployment} on a concrete case. Later renderings may refine or contest this entry; it remains as a
historical object.

% ============================================================================
\section{Feedback into the architecture}
% ============================================================================

Two FFGFT building blocks feed back directly into the governance: (i) the recoverability typology of Doc.~270
(three qualitatively different reduction/translation operations) as a concrete \emph{recoverability structure}
for the Matrix/OAP axis; (ii) the methodological guardrail of Doc.~190 (R35): every magnitude is writable as
$\xi^{N}$ or $\varphi^{N}$, a measured scale in an exponent is not a prediction -- only isolated, predicted,
precise \emph{ratios} carry forward content. Both operationalise Dot's admissibility and routing criteria on a
real example.

% ============================================================================
\section{Status}
% ============================================================================

A meta-level mapping, not a claim that either framework is correct; state June~2026, provisional and supersedable.
What it does: it shows, on principled grounds, \emph{where} the comparison from Doc.~271 is admissible
(spectral/recovery axis) and where not (question geometry, observer embedding). It remains open -- consistent with
Dot's non-absolutist stance -- whether the triadic resonance is more than coincidence and which substrate better
matches the observed mode spectrum.

% ============================================================================
\section{Sources}
% ============================================================================

\begin{itemize}
\item S.~Vossen, \emph{Dot Theory} (representational research programme) and the governance developed in the
direct IPI exchange (June~2026): \emph{Lexicon, OAP, Matrix, FAH} ($O0/O1/O2$ renderings; \enquote{declared
provenance, not universal agreement}; \enquote{preservation of meaning-making, not allocation of meaning}),
dottheory.co.uk.
\item M.~Krüger, HLV: \emph{Conservative Master Explanation} (Zenodo 10.5281/zenodo.20596861) and
\emph{6D-to-3D Cut-and-Project Graph Testbed} (Zenodo 10.5281/zenodo.20275888).
\item FFGFT: Doc.~190 (R35), Doc.~206 ($Z_3$), Doc.~270 (reduction chain / recoverability),
Doc.~271 (pairwise comparison FFGFT$\leftrightarrow$HLV).
\end{itemize}
